\begin{abstract}
Quantum reservoir computers process time series through the fading memory of a dissipative quantum system. We show
that the shape of the dissipation, set by its jump operators, matters as much as its rate. At equal total coupling,
spins that decay independently forget faster as the loss grows, whereas spins sharing one decay channel forget more
slowly: strong loss holds them in dark states through the quantum Zeno effect. A long memory has a price, because
later inputs scramble earlier ones into products that a linear readout cannot use. For a class of reservoirs we prove
that the number of inputs mixed into a stored memory is exponentially distributed, with mean equal to how hard one
input nudges it times how long it lasts. The input must therefore weaken as the memory grows, which more than doubles
the linear memory under exact readout. With finite sampling, weak local loss, which cools the spins, retains the most
memory; dark-state memories last as long but are over ten times fainter.
\end{abstract}
